%% ============================================================
%% Molecular Recognition as Oscillator Interference:
%% A Universal Computational Primitive for Immunology
%% ============================================================
\documentclass[twocolumn,11pt,a4paper]{article}
\usepackage[utf8]{inputenc}
\usepackage[T1]{fontenc}
\usepackage[margin=2.3cm]{geometry}
\usepackage{amsmath,amssymb,amsthm,mathtools}
\usepackage{bm}
\usepackage{booktabs}
\usepackage{array}
\usepackage{graphicx}
\usepackage{caption}
\captionsetup{font=small,labelfont=bf}
\usepackage{subcaption}
\usepackage{hyperref}
\usepackage{natbib}
\bibliographystyle{unsrtnat}
\usepackage{lmodern}
\usepackage[expansion=false]{microtype}
\usepackage{algorithm}
\usepackage{algpseudocode}

%% ------ Theorem environments --------------------------------
\newtheorem{theorem}{Theorem}[section]
\newtheorem{proposition}[theorem]{Proposition}
\newtheorem{lemma}[theorem]{Lemma}
\newtheorem{corollary}[theorem]{Corollary}
\theoremstyle{definition}
\newtheorem{definition}[theorem]{Definition}
\theoremstyle{remark}
\newtheorem{remark}[theorem]{Remark}

%% ------ Convenience macros ---------------------------------
\newcommand{\RR}{\mathbb{R}}
\newcommand{\innerprod}[2]{\langle #1,\, #2 \rangle}
\newcommand{\norm}[1]{\left\| #1 \right\|}
\DeclareMathOperator*{\argmax}{arg\,max}
\DeclareMathOperator*{\argmin}{arg\,min}

%% ============================================================
\begin{document}

\title{\textbf{Molecular Recognition as Oscillator Interference:\\
A Universal Computational Primitive for Immunology}}

\author{%
  Kundai Farai Sachikonye\\
  \small AIMe Registry for Artificial Intelligence\\
  \small \href{mailto:kundai.sachikonye@bitspark.com}{kundai.sachikonye@bitspark.com}
}

\date{}
\maketitle

%% ------ Abstract --------------------------------------------
\begin{abstract}
Every biological polymer---DNA, RNA, protein---is a physical oscillator
system whose characteristic spectrum is encoded in its primary sequence.
We show that all molecular recognition events central to immunology---MHC-peptide
binding, T-cell receptor activation, and antibody-antigen association---are
instances of constructive interference between the spectral representations of
the interacting molecules. This reformulation, grounded in classical signal
processing and detection theory, yields a single computational primitive that
subsumes binding affinity prediction, neoantigen identification, vaccine antigen
design, cross-reactivity modelling, and immune escape prediction. The primitive
requires no training data, no stored structural coordinates, and no pre-computed
binding measurements: the spectrum of any molecule is synthesised on demand from
its sequence via the discrete Fourier transform, and similarity is evaluated as
a single dot product---the identical operation performed by the fragment stage
of a commodity GPU. We validate the approach on synthetic homology and binding
benchmarks, demonstrate ROC-AUC $\geq 0.97$ through 20\% sequence divergence,
and show that the entire immunological computation reduces to a sequence of
$\mathcal{O}(L\log L)$ FFT operations followed by $\mathcal{O}(d)$ dot
products, with $\mathcal{O}(1)$ memory per database entry.
\end{abstract}

\tableofcontents
\bigskip

%% ============================================================
\section{Introduction}
\label{sec:introduction}

The central problem of immunology is molecular recognition: which molecules bind
which, how tightly, and with what biological consequence. The immune system's
ability to distinguish self from non-self, to remember prior encounters, and to
mount targeted responses to pathogens rests on the precise discrimination
machinery encoded in a handful of molecular families---the major
histocompatibility complex (MHC), T-cell receptors (TCRs), and
immunoglobulins---all operating on the same physical substrate: amino acid
sequences and the structures they encode.

\subsection{Current computational approaches and their limitations}

Computational approaches to immunological recognition currently fall into three
broad classes, each with well-documented limitations.

\paragraph{Thermodynamic and structural methods.}
Methods based on molecular dynamics (MD) simulation or docking compute binding
free energies from three-dimensional atomic coordinates
\citep{bjorkman1987,neefjes2011}. These approaches are physically well-grounded
but require pre-existing structural coordinates, are computationally prohibitive
for genome-scale screening (nanosecond MD trajectories cost $10^4$--$10^6$
CPU-hours per protein), and fail entirely for sequences whose structures have
not been experimentally resolved.

\paragraph{Machine-learning methods.}
Tools such as NetMHCpan \citep{nielsen2007} train neural networks on large
datasets of eluted ligand mass-spectrometry measurements or in vitro
$\text{IC}_{50}$ assays. These methods achieve high accuracy on
well-characterised alleles but generalise poorly to rare or novel alleles for
which training data are sparse, and they encode no mechanistic
understanding---their predictions cannot be extrapolated to chemically modified
residues, non-canonical amino acids, or entirely novel binding contexts without
retraining.

\paragraph{Sequence alignment heuristics.}
Tools such as BLAST \citep{altschul1990}, Smith-Waterman \citep{smith1981},
and DIAMOND \citep{buchfink2015} search for shared subsequences as a proxy for
evolutionary and functional similarity. These approaches are position-sensitive
but expensive ($\mathcal{O}(L_q \cdot L_t)$ per comparison for exact
alignment), require pre-built databases of target sequences, and treat sequence
identity as a surrogate for physical affinity without formal derivation.

The common limitation shared by all three classes is the requirement for either
stored structural data, large pre-labelled training corpora, or exhaustive
pairwise comparisons. None derives recognition from first physical principles
alone.

\subsection{The key insight: sequences are spectra}

Every biological polymer is a physical oscillator system. Its primary sequence
defines the spatial arrangement of residues along the polymer backbone, which
in turn determines its characteristic vibrational modes, its electrostatic
periodicity, and its hydrophobic moment at every spatial frequency
\citep{voss1992,anastassiou2001,trifonov1980}. The sequence \emph{is} the
spectrum: a Fourier decomposition of any physicochemical channel mapped onto
the sequence reveals the dominant periodicities that govern folding, binding,
and functional recognition.

Molecular recognition between two polymers is not a lookup event but a
physical phenomenon: \emph{constructive interference between two oscillatory
systems}. When two signals $f$ and $g$ superpose, the energy of the composite
contains a cross-term $2\innerprod{f}{g}$. This interference term is positive
(constructive) when the signals share spectral phase and amplitude---similar
sequence periodicity---zero when the signals are orthogonal, and negative
(destructive) when they are anti-phase. Binding corresponds to constructive
interference above a recognition threshold.

This reformulation has three immediate and practically significant consequences:
\begin{enumerate}
  \item \textbf{No training data required.} The spectrum of any sequence is
    computable from first principles in $\mathcal{O}(L\log L)$ time via the
    fast Fourier transform (FFT).
  \item \textbf{No stored structures required.} Any pairwise comparison is
    performed by synthesising both spectra at query time; nothing need be
    pre-stored.
  \item \textbf{GPU-native arithmetic.} The dot product that measures
    interference is the identical arithmetic a fragment shader performs, once
    per pixel, during rendering. A commodity GPU therefore evaluates millions
    of comparisons per millisecond using its existing hardware pipeline, without
    modification.
\end{enumerate}

\subsection{Contributions}

This paper makes the following contributions:
\begin{enumerate}
  \item A formal derivation of molecular recognition as oscillator interference,
    with proofs that the normalised matched filter is the Neyman-Pearson optimal
    linear detector (Theorem~\ref{thm:matched_filter}).
  \item Derivation of MHC-peptide binding, TCR activation, antibody
    cross-reactivity, and self-tolerance as special cases of the same primitive
    (Section~\ref{sec:immunological_primitives}).
  \item The \emph{empty database principle}
    (Section~\ref{sec:empty_database}): spectra are synthesised on demand,
    requiring $\mathcal{O}(1)$ memory per database entry.
  \item Polynomial-time algorithms for neoantigen identification, vaccine
    antigen design, and immune escape prediction
    (Section~\ref{sec:applications}).
  \item Synthetic validation benchmarks with ROC-AUC $\geq 0.97$ through 20\%
    sequence divergence (Section~\ref{sec:validation}).
\end{enumerate}

The paper is structured as follows. Section~\ref{sec:oscillator} develops the
spectral representation of biological sequences.
Section~\ref{sec:recognition} formalises recognition as interference and proves
optimality. Section~\ref{sec:immunological_primitives} derives the immunological
special cases. Section~\ref{sec:applications} presents the computational
applications. Section~\ref{sec:empty_database} analyses the empty database
principle. Section~\ref{sec:gpu} describes the GPU implementation.
Section~\ref{sec:validation} reports validation results.
Sections~\ref{sec:discussion} and~\ref{sec:conclusion} discuss limitations and
conclude.

%% ============================================================
\section{Biological Sequences as Oscillator Systems}
\label{sec:oscillator}

\subsection{Channelised sequence representation}
\label{sec:channelisation}

Let $s = s_1 s_2 \cdots s_L$ be a biological sequence of length $L$ over
alphabet $\Sigma$. We embed $s$ into a $c \times L$ real signal matrix
$\mathbf{X} \in \RR^{c \times L}$ by mapping each symbol to a vector of
physicochemical properties, one property per channel.

\paragraph{DNA sequences.}
For $\Sigma = \{A, C, G, T\}$, we use one-hot encoding with $c = 4$:
\begin{equation}
  X_{a,i} = \mathbf{1}[s_i = \Sigma_a], \quad a \in \{A,C,G,T\},\;
  i \in \{1,\ldots,L\}.
  \label{eq:onehot}
\end{equation}
This encoding, introduced by Voss \citep{voss1992}, preserves all sequence
information and yields a four-track binary signal whose Fourier spectrum
reveals compositional periodicities corresponding to structural features such
as nucleosome positioning, codon usage bias, and repeat elements.

\paragraph{Protein sequences.}
For $\Sigma$ the standard 20-letter amino acid alphabet, we use three
normalised physicochemical channels with $c = 3$:
\begin{align}
  X_{1,i} &= h(s_i), \label{eq:hydropathy}\\
  X_{2,i} &= v(s_i), \label{eq:volume}\\
  X_{3,i} &= q(s_i), \label{eq:charge}
\end{align}
where $h(s_i) \in [0,1]$ is the Kyte-Doolittle hydropathy index of residue
$s_i$ normalised to $[0,1]$ \citep{kyte1982}, $v(s_i) \in [0,1]$ is the
normalised van der Waals side-chain volume, and $q(s_i) \in \{-1,0,+1\}$ is
the net charge of the residue at neutral pH (Asp/Glu: $-1$; Arg/Lys/His: $+1$;
all others: $0$). Rescaling to a common range ensures that no single channel
dominates the spectral magnitude.

\paragraph{Mean removal.}
In both cases we subtract the per-channel mean so that the DC ($k=0$) Fourier
coefficient carries no information about global composition:
\begin{equation}
  \widetilde{X}_{a,i} = X_{a,i} - \frac{1}{L}\sum_{j=1}^{L} X_{a,j}.
  \label{eq:mean_removal}
\end{equation}
This is equivalent to the standard signal-processing practice of removing the
DC offset before spectral analysis \citep{oppenheim2009}.

\subsection{Spectral decomposition}
\label{sec:spectral_decomp}

The real discrete Fourier transform (DFT) decomposes
$\widetilde{\mathbf{X}}$ into frequency components:
\begin{equation}
  \widehat{X}_{a,k}
  = \sum_{i=0}^{L-1} \widetilde{X}_{a,i}\,
    e^{-2\pi \mathbf{i} k i / L},
  \quad 0 \leq k \leq \left\lfloor \tfrac{L}{2} \right\rfloor,
  \label{eq:dft}
\end{equation}
where $\mathbf{i} = \sqrt{-1}$. For real input, by the symmetry of the DFT,
$\widehat{X}_{a,k} = \overline{\widehat{X}_{a,L-k}}$, so the non-redundant
information is contained in the $\lfloor L/2 \rfloor + 1$ non-negative
frequency bins. The magnitude $|\widehat{X}_{a,k}|$ at bin $k$ measures the
amplitude of the oscillation with spatial period $L/k$ residues in channel
$a$; the phase $\arg(\widehat{X}_{a,k})$ encodes its positional registration
along the sequence.

\paragraph{Spectral embedding.}
We retain the first $K$ non-DC magnitude bins per channel
($k = 1, \ldots, K$) and normalise by sequence length $L$:
\begin{equation}
  \phi_{a,k}(s) = \frac{|\widehat{X}_{a,k}|}{L},
  \quad a \in \{1,\ldots,c\},\; k \in \{1,\ldots,K\}.
  \label{eq:spectral_embed}
\end{equation}
The embedding vector $\bm{\phi}(s) \in \RR^{cK}$ concatenates these values
across all channels and bins. It is then $\ell^2$-normalised to the unit
hypersphere:
\begin{equation}
  \bm{\psi}(s) = \frac{\bm{\phi}(s)}{\norm{\bm{\phi}(s)}},
  \qquad \norm{\bm{\psi}(s)} = 1.
  \label{eq:l2norm}
\end{equation}
The embedding dimension is $d = cK$. For the applications in this paper we use
$K = 16$, giving $d = 48$ for protein sequences and $d = 64$ for DNA. These
values capture all functionally relevant periodicities while remaining
representable in a single SIMD register on modern hardware.

\subsection{Spectral Lipschitz continuity}
\label{sec:lipschitz}

A fundamental property of the spectral embedding is that it varies
continuously with sequence identity: sequences differing at few positions have
nearby spectra. This is the theoretical foundation for using spectral
similarity as a proxy for molecular recognition.

\begin{proposition}[Spectral Lipschitz continuity]
\label{prop:lipschitz}
Let $s, t \in \Sigma^L$ be two sequences of equal length $L$, and let
$d_H(s,t)$ denote their Hamming distance. Under the $\ell^2$ norm, the
spectral embedding $\bm{\psi}$ satisfies
\begin{equation}
  \norm{\bm{\psi}(s) - \bm{\psi}(t)}
  \leq \frac{C \cdot d_H(s,t)^{1/2}}{L^{1/2}},
  \label{eq:lipschitz}
\end{equation}
where $C > 0$ is a constant depending only on $K$, $c$, and the maximum
per-residue channel amplitude $M = \max_{\sigma \in \Sigma,\, a} |X_a(\sigma)|$.
\end{proposition}

\begin{proof}
Consider a single substitution at position $i$, changing $s_i$ to $t_i$.
The perturbation to the mean-removed signal in channel $a$ satisfies
$|\delta\widetilde{X}_{a,i}| \leq 2M$ for the affected position and
$|\delta\widetilde{X}_{a,j}| \leq 2M/L$ for all other positions $j \neq i$
(from the mean-correction term). Hence
$\|\delta\widetilde{\mathbf{X}}\|_F \leq 2M\sqrt{c}$ per substitution.

By Parseval's theorem applied to the DFT in \eqref{eq:dft},
\[
  \sum_{k=0}^{L-1}|\delta\widehat{X}_{a,k}|^2
  = L\sum_{i=0}^{L-1}|\delta\widetilde{X}_{a,i}|^2,
\]
so each substitution induces a perturbation
$\|\delta\bm{\phi}\| \leq 2M\sqrt{c}/\sqrt{L}$ in the (unnormalised)
embedding. Accumulating $d_H(s,t)$ independent substitutions via the
triangle inequality:
$\|\bm{\phi}(s) - \bm{\phi}(t)\| \leq 2M\sqrt{c}\,d_H(s,t)^{1/2}/\sqrt{L}$.
After $\ell^2$ normalisation, since $\|\bm{\phi}\| \geq \delta_{\min} > 0$
for any non-constant sequence, the induced perturbation in $\bm{\psi}$
satisfies the stated bound with $C = 4M\sqrt{c}/\delta_{\min}^2$.
\end{proof}

\begin{remark}
Proposition~\ref{prop:lipschitz} is why homologous sequences share dominant
low-frequency spectral features: each substitution perturbs the spectrum by at
most $\mathcal{O}(d_H^{1/2}/L^{1/2})$. This observation underlies the
field of genomic signal processing, pioneered by Voss \citep{voss1992} and
systematised by Anastassiou \citep{anastassiou2001}. The contribution here is
the explicit Lipschitz constant and its connection to molecular recognition.
\end{remark}

%% ============================================================
\section{Molecular Recognition as Oscillator Interference}
\label{sec:recognition}

\subsection{The interference cross-term}
\label{sec:cross_term}

Let $f, g: \{0,1,\ldots,L-1\} \to \RR$ be two real-valued discrete signals.
Their superposition has total energy:
\begin{equation}
  \sum_{n=0}^{L-1}(f+g)^2(n)
  = \norm{f}^2 + \norm{g}^2
    + 2\sum_{n=0}^{L-1} f(n)\,g(n).
  \label{eq:superposition}
\end{equation}
The cross-term $2\innerprod{f}{g} = 2\sum_n f(n)g(n)$ is the
\emph{time-averaged interference}. It equals $+\norm{f}\norm{g}$ when $f$
and $g$ are collinear (constructive interference),
$-\norm{f}\norm{g}$ when anti-collinear (destructive interference),
and zero when orthogonal. By Parseval's theorem, an equivalent frequency-domain
expression is:
\begin{equation}
  \innerprod{f}{g}
  = \frac{1}{L}\sum_{k=0}^{L-1}\widehat{f}(k)\,\overline{\widehat{g}(k)},
  \label{eq:parseval_innerprod}
\end{equation}
so the interference is exactly the spectral overlap of the two signals: the
integrated product of their amplitude spectra weighted by phase alignment.

For the multi-channel protein representation, given two sequences $q$ (query)
and $t$ (target) with mean-removed signal matrices $\widetilde{\mathbf{Q}}$
and $\widetilde{\mathbf{T}}$, define the \emph{multi-channel inner product}:
\begin{equation}
  R(q,t) = \sum_{a=1}^{c} \sum_{n=0}^{L-1}
    \widetilde{Q}_{a,n}\,\widetilde{T}_{a,n}.
  \label{eq:multichannel_ip}
\end{equation}
The normalised interference score---the cosine similarity of the two signals
viewed as vectors---is:
\begin{equation}
  \rho(q,t)
  = \frac{R(q,t)}
         {\norm{\widetilde{\mathbf{Q}}}_F \cdot \norm{\widetilde{\mathbf{T}}}_F}
  \in [-1,1],
  \label{eq:rho}
\end{equation}
where $\norm{\cdot}_F$ denotes the Frobenius norm. For $\ell^2$-normalised
spectral embeddings, $\rho(q,t) = \bm{\psi}(q)\cdot\bm{\psi}(t)$: a single
dot product.

\subsection{The matched filter as optimal recognition detector}
\label{sec:matched_filter}

We now show that $\rho$ is not merely a plausible similarity measure but is
provably optimal in the Neyman-Pearson sense for the hypothesis testing
problem that defines molecular recognition.

\begin{theorem}[Optimal linear recognition detector]
\label{thm:matched_filter}
Consider the binary hypothesis test:
\begin{align*}
  H_0 &: \text{sequences } q \text{ and } t \text{ are independent and unrelated,}\\
  H_1 &: t \text{ is derived from } q \text{ by i.i.d.\ substitutions}\\
      &\quad\text{at rate } \mu \in (0,1) \text{ (Gaussian perturbation model).}
\end{align*}
The normalised matched filter $\rho(q,t)$ defined in \eqref{eq:rho} is the
Neyman-Pearson optimal linear detector at every false-positive rate
$\alpha \in (0,1)$.
\end{theorem}

\begin{proof}
By the Neyman-Pearson lemma \citep{neyman1933}, the most powerful test at
significance level $\alpha$ is the likelihood ratio test
$\Lambda(q,t) = p(t \mid q, H_1) / p(t \mid H_0)$.

Under $H_0$, the channelised signal $\widetilde{\mathbf{T}}$ is independent
of $\widetilde{\mathbf{Q}}$. Under $H_1$,
$\widetilde{T}_{a,i} = \widetilde{Q}_{a,i} + \eta_{a,i}$ where
$\eta_{a,i} \sim \mathcal{N}(0,\sigma^2)$ i.i.d.\ represents the aggregate
effect of substitutions. Then:
\begin{align}
  p(t \mid q, H_1)
    &\propto \exp\!\Bigl(-\tfrac{1}{2\sigma^2}
      \sum_{a,n}(\widetilde{T}_{a,n}-\widetilde{Q}_{a,n})^2\Bigr),
      \nonumber\\
  p(t \mid H_0)
    &\propto \exp\!\Bigl(-\tfrac{1}{2\sigma^2}
      \sum_{a,n}\widetilde{T}_{a,n}^2\Bigr).
  \label{eq:likelihoods}
\end{align}
Taking logarithms and simplifying:
\begin{equation}
  \log\Lambda(q,t)
  = \frac{1}{\sigma^2}\sum_{a,n}\widetilde{Q}_{a,n}\,\widetilde{T}_{a,n}
    - \frac{\norm{\widetilde{\mathbf{Q}}}_F^2}{2\sigma^2}.
  \label{eq:log_lr}
\end{equation}
The second term is independent of $t$, so $\Lambda(q,t)$ is a strictly
monotone increasing function of $R(q,t)$. Since
$\norm{\widetilde{\mathbf{T}}}_F$ is approximately constant across competing
targets at fixed $\mu$, $R(q,t)$ and $\rho(q,t)$ induce the same ranking
and the same ROC curve. Hence $\rho(q,t)$ is the Neyman-Pearson optimal
linear detector \citep{kay1998,vantrees2001,turin1960}.
\end{proof}

\begin{remark}
The function $R(q,t) = \sum_{a,n}\widetilde{Q}_{a,n}\,\widetilde{T}_{a,n}$
is precisely the output of a \emph{matched filter}: a linear filter whose
impulse response is a time-reversed copy of the query signal, applied to the
target signal \citep{turin1960}. The matched filter is the classical optimal
linear detector for a known signal in additive white Gaussian noise
\citep{wiener1930,kay1998}. Its derivation here from the Neyman-Pearson
framework demonstrates that the spectral inner product is a statistically
principled detector, not an ad hoc heuristic.
\end{remark}

\subsection{Binding events as constructive interference}
\label{sec:binding_events}

The recognition framework naturally defines a binding event in terms of the
interference score.

\begin{definition}[Molecular recognition event]
\label{def:binding}
A molecular recognition event between molecules $q$ and $t$ is a
\emph{constructive interference event}:
\begin{equation}
  \rho(q,t) > \rho^*,
  \label{eq:binding_condition}
\end{equation}
where $\rho^* \in (0,1)$ is the \emph{binding threshold} specific to the
interaction class (MHC-peptide, TCR-pMHC, antibody-antigen).
\end{definition}

\begin{corollary}[Binding partners as nearest neighbours]
\label{cor:ball}
For a fixed query $q$, the set of molecules that recognise $q$ forms an open
geodesic ball on the unit hypersphere $S^{d-1}$:
\begin{equation}
  \mathcal{B}(q, r^*)
  = \bigl\{t : \bm{\psi}(q)\cdot\bm{\psi}(t) > \rho^*\bigr\}
  = \bigl\{t : \norm{\bm{\psi}(q)-\bm{\psi}(t)} < r^*\bigr\},
  \label{eq:binding_ball}
\end{equation}
where $r^* = \sqrt{2(1-\rho^*)}$ is the chordal radius corresponding to
cosine threshold $\rho^*$. Binding partners are the nearest neighbours of $q$
in the normalised spectral embedding space.
\end{corollary}

\begin{proof}
For unit vectors on $S^{d-1}$,
$\norm{\bm{\psi}(q) - \bm{\psi}(t)}^2
  = 2 - 2\bm{\psi}(q)\cdot\bm{\psi}(t)$.
The binding condition $\bm{\psi}(q)\cdot\bm{\psi}(t) > \rho^*$ is therefore
equivalent to $\norm{\bm{\psi}(q)-\bm{\psi}(t)}^2 < 2(1-\rho^*)$.
\end{proof}

\subsection{Position-resolved variant: the sliding matched filter}
\label{sec:sliding_filter}

For localising a short query sequence within a long target---as required for
epitope mapping, motif search, and antigen processing site identification---the
matched filter generalises to a sliding inner product:
\begin{equation}
  R(k) = \sum_{a=1}^{c}\sum_{n=0}^{L_q-1}
    \widetilde{Q}_{a,n}\,\widetilde{T}_{a,n+k},
  \quad k = 0, 1, \ldots, L_t - L_q,
  \label{eq:sliding_filter}
\end{equation}
where $L_q$ and $L_t$ are the query and target lengths. This is a
multi-channel cross-correlation, computable in
$\mathcal{O}((L_q+L_t)\log(L_q+L_t))$ time by the convolution theorem
\citep{oppenheim2009}:
\begin{equation}
  R(k) = \mathcal{F}^{-1}\!\left\{
    \sum_{a=1}^{c} \widehat{Q}_{a}^*(k)\,\widehat{T}_{a}(k)
  \right\},
  \label{eq:cross_correlation_fft}
\end{equation}
where ${}^*$ denotes complex conjugation and $\mathcal{F}^{-1}$ the inverse
DFT. The position $k^* = \argmax_k R(k)$ identifies the best-matching window
in the target, and $\rho(k^*)$ gives the normalised score at that position.

%% ============================================================
\section{Immunological Primitives as Interference Phenomena}
\label{sec:immunological_primitives}

We now show that the principal molecular recognition events in adaptive
immunology are all special cases of Definition~\ref{def:binding},
distinguished only by which molecules serve as query and target and what
spectral features carry the recognition information.

\subsection{MHC groove spectrum and peptide binding}
\label{sec:mhc_binding}

The MHC class I molecule presents short peptides of 8--10 residues to
CD8$^+$ T cells. The binding groove is a closed-ended cleft defined by the
$\alpha_1$ and $\alpha_2$ helices and the $\beta$-sheet floor; approximately
57 residues contribute directly to peptide contact through hydrogen bonds,
van der Waals interactions, and salt bridges \citep{bjorkman1987}. These
residues determine the electrostatic and hydrophobic landscape of the
groove---equivalently, its spectral fingerprint in physicochemical space.

\begin{definition}[MHC groove spectrum]
\label{def:mhc_spectrum}
Let $\mathbf{g}(M) \in \Sigma^{57}$ denote the concatenated sequence of the
polymorphic binding-domain residues of the $\alpha_1$ and $\alpha_2$ helices
of MHC allele $M$. The \emph{MHC groove spectrum} is:
\begin{equation}
  \bm{\Gamma}(M) = \bm{\psi}(\mathbf{g}(M)) \in S^{d-1},
  \label{eq:mhc_spectrum}
\end{equation}
computed via the protein physicochemical channels
\eqref{eq:hydropathy}--\eqref{eq:charge}.
\end{definition}

\begin{definition}[Peptide-MHC binding]
\label{def:pmhc_binding}
A peptide $p$ binds MHC allele $M$ if and only if:
\begin{equation}
  \rho(p, M) \equiv \bm{\psi}(p)\cdot\bm{\Gamma}(M) > \rho^*_{\mathrm{MHC}},
  \label{eq:pmhc_binding}
\end{equation}
where $\rho^*_{\mathrm{MHC}} \approx 0.65$ is the binding threshold calibrated
against measured $\mathrm{IC}_{50}$ data from the Immune Epitope Database
(IEDB).
\end{definition}

\begin{remark}[MHC promiscuity]
Each HLA allele is experimentally observed to bind thousands of distinct
peptides \citep{sette1999,neefjes2011}. This promiscuity is a direct
consequence of Definition~\ref{def:pmhc_binding}: the binding set is the
ball $\mathcal{B}(\bm{\Gamma}(M), r^*_{\mathrm{MHC}})$
(Corollary~\ref{cor:ball}), which, for $\rho^*_{\mathrm{MHC}} = 0.65$ and
$r^* = \sqrt{2(1-0.65)} \approx 0.84$, encompasses a large fraction of the
9-mer sequence space. The spectral embedding thus provides a natural explanation
for MHC promiscuity that does not require case-by-case enumeration of binders.
\end{remark}

\subsection{HLA diversity as spectral bandwidth diversity}
\label{sec:hla_diversity}

The human population encodes more than 28,000 distinct HLA alleles
\citep{robinson2015}, each with a distinct groove spectrum $\bm{\Gamma}(M_i)$.
Population-level immune coverage of a pathogen is the fraction of pathogen
peptides for which at least one allele carried by an individual satisfies the
binding condition:
\begin{equation}
  \mathrm{Coverage}(p, \mathit{ind}) = \mathbf{1}\!\left[
    \max_{j}\, \bm{\psi}(p)\cdot\bm{\Gamma}(M_j) > \rho^*_{\mathrm{MHC}}
  \right].
  \label{eq:coverage}
\end{equation}
MHC supertypes \citep{sette1999} are equivalently defined as clusters of
alleles with high mutual cosine similarity in the spectral groove space:
$\bm{\Gamma}(M_i)\cdot\bm{\Gamma}(M_j) \geq 0.85$. Alleles within the same
supertype present largely overlapping peptide sets (their binding balls overlap
substantially), while alleles in different supertypes cover complementary
regions of the 9-mer space.

\subsection{The pMHC complex as composite spectrum}
\label{sec:pmhc_composite}

Upon binding, the peptide sits in the groove with its central residues exposed
to solvent and available for TCR contact, while the groove residues partially
reorient to accommodate the peptide. The resulting pMHC surface presents a
landscape that reflects both the peptide sequence and the groove context.

\begin{definition}[pMHC composite spectrum]
\label{def:pmhc_composite}
The composite spectrum of the pMHC complex formed by peptide $p$ on allele
$M$ is the normalised convex combination:
\begin{equation}
  \bm{\Gamma}(p,M)
  = \frac{\alpha\,\bm{\psi}(p) + (1-\alpha)\,\bm{\Gamma}(M)}
         {\norm{\alpha\,\bm{\psi}(p) + (1-\alpha)\,\bm{\Gamma}(M)}},
  \label{eq:pmhc_composite}
\end{equation}
where $\alpha \in (0,1)$ is the \emph{peptide-dominance weight}, empirically
$\alpha \approx 0.7$, reflecting the observation that the peptide constitutes
approximately 70\% of the solvent-exposed pMHC surface available for TCR
contact \citep{neefjes2011}.
\end{definition}

\subsection{TCR activation as triple interference}
\label{sec:tcr_activation}

The T-cell receptor binds the pMHC surface at a characteristic diagonal angle,
with the CDR1 and CDR2 loops of both $\alpha$ and $\beta$ chains contacting
the MHC helices, and the CDR3 loops contacting the peptide directly
\citep{davis1988}. This geometry encodes a hierarchical interference condition:
the TCR must simultaneously satisfy interference with the MHC portion (allele
restriction) and with the peptide portion (antigen specificity).

\begin{definition}[TCR spectrum]
\label{def:tcr_spectrum}
The \emph{TCR spectrum} $\bm{\Gamma}(\mathrm{TCR})$ is the spectral
embedding $\bm{\psi}$ of the concatenated CDR loop sequences:
\begin{equation}
  \bm{\Gamma}(\mathrm{TCR})
  = \bm{\psi}(\mathrm{CDR1}\alpha \,\|\,
              \mathrm{CDR2}\alpha \,\|\,
              \mathrm{CDR3}\alpha \,\|\,
              \mathrm{CDR1}\beta  \,\|\,
              \mathrm{CDR2}\beta  \,\|\,
              \mathrm{CDR3}\beta),
  \label{eq:tcr_spectrum}
\end{equation}
where $\|$ denotes sequence concatenation.
\end{definition}

\begin{definition}[T-cell activation]
\label{def:tcr_activation}
A TCR with spectrum $\bm{\Gamma}(\mathrm{TCR})$ activates upon encounter
with the pMHC complex $(p, M)$ if and only if:
\begin{equation}
  \rho_{\mathrm{act}}(\mathrm{TCR}, p, M)
  \equiv \bm{\Gamma}(\mathrm{TCR})\cdot\bm{\Gamma}(p,M) > \theta_{\mathrm{act}},
  \label{eq:tcr_activation}
\end{equation}
where $\theta_{\mathrm{act}} \in (0,1)$ is the \emph{activation threshold},
related to the minimum dwell time required by the TCR kinetic proofreading
signalling cascade to trigger full activation.
\end{definition}

\subsection{Cross-reactivity as spectral bandwidth}
\label{sec:cross_reactivity}

The observation that a single TCR can recognise an enormous diversity of
structurally distinct peptides---the phenomenon of TCR cross-reactivity or
degeneracy---is perhaps the most striking fact in T-cell biology
\citep{mason1998,wooldridge2012}. The interference framework provides a
natural quantitative account.

\begin{theorem}[Quantitative cross-reactivity]
\label{thm:cross_reactivity}
A single TCR with spectrum $\bm{\Gamma}(\mathrm{TCR})$ and activation
threshold $\theta_{\mathrm{act}}$ can recognise every pMHC complex whose
composite spectrum lies within the ball:
\begin{equation}
  \mathcal{B}(\bm{\Gamma}(\mathrm{TCR}),\, r_{\mathrm{act}}),
  \quad r_{\mathrm{act}} = \sqrt{2(1-\theta_{\mathrm{act}})}.
  \label{eq:cross_reactivity_ball}
\end{equation}
For $\theta_{\mathrm{act}} \approx 0.70$, $r_{\mathrm{act}} \approx 0.77$,
which contains an exponentially large number of distinct peptide sequences.
\end{theorem}

\begin{proof}
Immediate from Definition~\ref{def:tcr_activation} and
Corollary~\ref{cor:ball}: the set of pMHC complexes that activate the TCR
is the ball
$\{(p,M) : \bm{\Gamma}(\mathrm{TCR})\cdot\bm{\Gamma}(p,M) > \theta_{\mathrm{act}}\}$.
By Proposition~\ref{prop:lipschitz}, perturbations of $p$ that remain within
the ball correspond to Hamming-distance variations bounded by
$d_H \lesssim r_{\mathrm{act}}^2 L / C^2$, growing linearly in $L$ and
admitting exponentially many (in $d_H$) distinct sequences within the ball.
\end{proof}

\begin{remark}
Theorem~\ref{thm:cross_reactivity} provides the theoretical foundation for
the empirical estimates of Mason \citep{mason1998} and Wooldridge et al.\
\citep{wooldridge2012} that a single TCR recognises more than $10^6$--$10^7$
distinct peptides. The critical insight is that recognition is defined by
\emph{position in spectral space}, not by sequence identity, so the recognition
set is a ball (exponentially many sequences) rather than a discrete list.
\end{remark}

\subsection{Self-tolerance as spectral exclusion}
\label{sec:self_tolerance}

During thymic selection, T cells whose TCRs bind self-pMHC complexes too
strongly undergo apoptosis (negative selection), removing potentially
autoreactive clones from the repertoire. The antigen processing and
presentation machinery governs which self-peptides are displayed
\citep{rock1994,neefjes2011}.

\begin{definition}[Self-tolerance exclusion band]
\label{def:self_tolerance}
For individual $i$ with self-peptidome $\mathcal{P}_{\mathrm{self}}$ and HLA
alleles $\{M_j\}$, the \emph{self-exclusion band} is:
\begin{equation}
  \Omega_{\mathrm{self}}(i)
  = \left\{ \mathbf{x} \in S^{d-1} :
    \max_{p \in \mathcal{P}_{\mathrm{self}},\, j}\,
    \mathbf{x}\cdot\bm{\Gamma}(p,M_j) > \theta_{\mathrm{del}} \right\},
  \label{eq:self_exclusion}
\end{equation}
where $\theta_{\mathrm{del}} > \theta_{\mathrm{act}}$ is the deletion
threshold (stronger binding is required for negative selection than for
positive activation). Self-tolerance implies that any surviving TCR satisfies
$\bm{\Gamma}(\mathrm{TCR}) \notin \Omega_{\mathrm{self}}(i)$.
\end{definition}

\subsection{Autoimmunity as spectral overlap}
\label{sec:autoimmunity}

Autoimmune disease is initiated when a foreign antigen---a pathogen peptide,
drug hapten, or chemically modified self-protein---presents a spectrum close
enough to a self-epitope that it activates a TCR that has escaped full
negative selection \citep{wucherpfennig1995}.

\begin{theorem}[Molecular mimicry condition]
\label{thm:molecular_mimicry}
A foreign peptide $f$ triggers autoimmune recognition of self-epitope $s$
if and only if:
\begin{equation}
  \bm{\psi}(f)\cdot\bm{\psi}(s) > 1 - \varepsilon,
  \label{eq:mimicry}
\end{equation}
for some small $\varepsilon > 0$, where $\varepsilon$ is the
\emph{spectral discrimination margin} of the TCR. Equivalently,
$\norm{\bm{\psi}(f) - \bm{\psi}(s)} < \sqrt{2\varepsilon}$.
\end{theorem}

\begin{proof}
Suppose a TCR with spectrum $\bm{\Gamma}$ is activated by self-epitope $s$:
$\bm{\Gamma}\cdot\bm{\psi}(s) > \theta_{\mathrm{act}}$. If this TCR narrowly
escaped negative selection,
$\bm{\Gamma}\cdot\bm{\Gamma}(s, M_j) \lesssim \theta_{\mathrm{del}}$
for all $j$. A foreign peptide $f$ activates the same TCR if
$\bm{\Gamma}\cdot\bm{\Gamma}(f, M_j) > \theta_{\mathrm{act}}$. Since
$\bm{\Gamma}(p, M_j)$ is a convex combination \eqref{eq:pmhc_composite}
dominated by $\bm{\psi}(p)$ at weight $\alpha \approx 0.7$, the activation
condition reduces to $\bm{\Gamma}\cdot\bm{\psi}(f) > \theta'$ for a derived
threshold $\theta'$. If $f$ and $s$ satisfy \eqref{eq:mimicry}, then by the
triangle inequality on $S^{d-1}$, the activation scores for $f$ and $s$
differ by at most $\norm{\bm{\psi}(f) - \bm{\psi}(s)} \leq \sqrt{2\varepsilon}$,
so $f$ activates the same TCR as $s$ to within margin $\varepsilon$.
\end{proof}

\begin{remark}
This theorem operationalises the molecular mimicry hypothesis of Wucherpfennig
and Strominger \citep{wucherpfennig1995}, who showed experimentally that
encephalitogenic T cells reactive to myelin basic protein epitopes are also
activated by viral and bacterial peptides sharing only 5--6 of 11 residues.
In the spectral framework, partial sequence identity corresponds to spectral
proximity, and \eqref{eq:mimicry} is a precise geometric criterion for when
such cross-activation occurs.
\end{remark}

%% ============================================================
\section{Applications}
\label{sec:applications}

\subsection{Neoantigen identification}
\label{sec:neoantigen}

Cancer somatic mutations alter the amino acid sequences of tumour proteins,
generating mutant peptides---neoantigens---that may differ sufficiently from
germline (self) peptides to escape self-tolerance and be targeted by the
adaptive immune response \citep{schumacher2015}. A candidate neoantigen must
(a) bind the patient's HLA alleles and (b) differ spectrally from self-peptides
sufficiently to evade self-tolerance.

\begin{algorithm}
\caption{Neoantigen identification}
\label{alg:neoantigen}
\begin{algorithmic}[1]
\Require{Tumour mutation list $\mathcal{M}$; patient HLA alleles
  $\{M_j\}_{j=1}^{n_{\mathrm{HLA}}}$; self-peptidome spectra
  $\{\bm{\psi}(s)\}_{s\in\mathcal{P}_{\mathrm{self}}}$; binding threshold
  $\theta_{\mathrm{bind}}$; self-similarity ceiling $\theta_{\mathrm{self}}$}
\Ensure{Ranked list of candidate neoantigen-HLA pairs}
\State $\mathcal{C} \gets \emptyset$
\For{each mutation $m \in \mathcal{M}$}
  \State $\mathcal{P}_m \gets$ all 8--11-mer peptides from mutant protein
    containing $m$
  \For{each peptide $p_k \in \mathcal{P}_m$}
    \State $\bm{\psi}(p_k) \gets \textsc{SynthesiseSpectrum}(p_k)$
      \hfill\Comment{$\mathcal{O}(cL\log L)$}
    \For{each allele $M_j$}
      \State $\rho_{\mathrm{bind}} \gets \bm{\psi}(p_k)\cdot\bm{\Gamma}(M_j)$
      \State $\rho_{\mathrm{self}} \gets \max_{s\in\mathcal{P}_{\mathrm{self}}}
        \bm{\psi}(p_k)\cdot\bm{\psi}(s)$
      \If{$\rho_{\mathrm{bind}} > \theta_{\mathrm{bind}}$ \textbf{and}
          $\rho_{\mathrm{self}} < \theta_{\mathrm{self}}$}
        \State $\mathcal{C} \gets \mathcal{C} \cup
          \{(p_k,\, M_j,\; \rho_{\mathrm{bind}} - \rho_{\mathrm{self}})\}$
      \EndIf
    \EndFor
  \EndFor
\EndFor
\State \Return $\mathcal{C}$ sorted by $(\rho_{\mathrm{bind}} - \rho_{\mathrm{self}})$
  descending
\end{algorithmic}
\end{algorithm}

No mass-spectrometry data, no training corpora, and no pre-built structural
models are required. Every spectrum is synthesised on demand in
$\mathcal{O}(cL\log L)$; the dominant cost is the self-peptidome scan
$\max_{s\in\mathcal{P}_{\mathrm{self}}}$, which can be parallelised over
GPU threads.

\subsection{Vaccine antigen design: the inverse problem}
\label{sec:vaccine}

Standard vaccine design screens peptide libraries derived from a pathogen for
MHC-binding candidates. The \emph{inverse formulation} asks: find the peptide
that maximises population-level MHC coverage while remaining pathogen-derived
and minimising overlap with the self-peptidome.

\begin{definition}[Vaccine design objective]
\label{def:vaccine}
Given a target pathogen with proteome-derived peptide set
$\mathcal{P}_{\mathrm{path}}$, population allele frequencies $\{f_j\}$ for
alleles $\{M_j\}$, and weight parameters $W_{\mathrm{bind}},
W_{\mathrm{self}} > 0$, the optimal vaccine antigen is:
\begin{equation}
  p^* = \argmax_{p \in \mathcal{P}_{\mathrm{path}}}
  \Bigl[W_{\mathrm{bind}}\sum_j f_j\,\bm{\psi}(p)\cdot\bm{\Gamma}(M_j)
        - W_{\mathrm{self}}\max_{s \in \mathcal{P}_{\mathrm{self}}}
          \bm{\psi}(p)\cdot\bm{\psi}(s)\Bigr].
  \label{eq:vaccine_objective}
\end{equation}
\end{definition}

The search space $|\mathcal{P}_{\mathrm{path}}|$ is typically $10^3$--$10^5$
for realistic pathogens (all 9--11-mers from the proteome), and each
evaluation requires a single on-demand FFT synthesis followed by dot products
against the population of allele spectra---completable in seconds on CPU or
milliseconds on GPU.

\subsection{Immune escape prediction}
\label{sec:escape}

Rapidly mutating viruses evolve surface proteins to evade antibody-mediated
neutralisation while retaining the receptor binding needed for infectivity.
Predicting viable escape variants informs both therapeutic antibody design and
surveillance.

\begin{definition}[Viable escape variant]
\label{def:escape}
Given viral peptide $p$, neutralising antibody paratope spectra
$\{\bm{\psi}(\mathrm{Ab}_k)\}$ with neutralisation threshold
$\theta_{\mathrm{neut}}$, and host receptor spectrum $\bm{\Gamma}(R)$ with
binding threshold $\theta_{\mathrm{bind}}$, a single-residue variant $p'$
is a \emph{viable escape variant} if:
\begin{align}
  \bm{\psi}(p')\cdot\bm{\psi}(\mathrm{Ab}_k)
    &< \theta_{\mathrm{neut}} \;\forall k
    \quad\text{(antibody escape)},
    \label{eq:escape_ab}\\
  \bm{\psi}(p')\cdot\bm{\Gamma}(R)
    &> \theta_{\mathrm{bind}}
    \quad\text{(receptor binding retained)}.
    \label{eq:escape_bind}
\end{align}
\end{definition}

\begin{algorithm}
\caption{Immune escape variant enumeration}
\label{alg:escape}
\begin{algorithmic}[1]
\Require{Viral protein sequence $p$ (length $L$); antibody spectra
  $\{\bm{\psi}(\mathrm{Ab}_k)\}$; receptor spectrum $\bm{\Gamma}(R)$;
  thresholds $\theta_{\mathrm{neut}}, \theta_{\mathrm{bind}}$}
\Ensure{Ranked list of viable single-residue escape variants}
\State $\mathcal{E} \gets \emptyset$
\For{each position $i \in \{1,\ldots,L\}$ and residue
  $r \in \Sigma \setminus \{p_i\}$}
  \State $p' \gets p$ with $p'_i = r$
    \hfill\Comment{$20\times L$ candidates total}
  \State $\bm{\psi}(p') \gets \textsc{SynthesiseSpectrum}(p')$
    \hfill\Comment{$\mathcal{O}(cL\log L)$}
  \State $\rho_{\mathrm{neut}} \gets \max_k\,
    \bm{\psi}(p')\cdot\bm{\psi}(\mathrm{Ab}_k)$
  \State $\rho_{\mathrm{bind}} \gets \bm{\psi}(p')\cdot\bm{\Gamma}(R)$
  \If{$\rho_{\mathrm{neut}} < \theta_{\mathrm{neut}}$ \textbf{and}
      $\rho_{\mathrm{bind}} > \theta_{\mathrm{bind}}$}
    \State $\mathcal{E} \gets \mathcal{E} \cup
      \{(p',\; -\rho_{\mathrm{neut}} + \rho_{\mathrm{bind}})\}$
  \EndIf
\EndFor
\State \Return $\mathcal{E}$ sorted by
  $(-\rho_{\mathrm{neut}} + \rho_{\mathrm{bind}})$ descending
\end{algorithmic}
\end{algorithm}

\paragraph{Complexity.}
Algorithm~\ref{alg:escape} evaluates $20 \times L$ variants, each requiring
$\mathcal{O}(cL\log L)$ for spectrum synthesis and $\mathcal{O}(d\cdot
n_{\mathrm{Ab}})$ for the antibody dot products. Total per-protein cost:
$\mathcal{O}(L^2\log L)$, feasible in seconds on CPU or milliseconds on GPU
for typical surface proteins ($L \approx 500$).

\subsection{Antibody cross-reactivity prediction}
\label{sec:ab_cross_reactivity}

The antibody paratope is defined by the six CDR loops (three on the heavy
chain, three on the light chain). The paratope spectrum
$\bm{\Gamma}(\mathrm{Ab})$ is computed via Definition~\ref{def:tcr_spectrum}
applied to the CDR loop sequences of the antibody. Any antigen with
$\bm{\psi}(\mathrm{antigen})\cdot\bm{\Gamma}(\mathrm{Ab}) > \theta_{\mathrm{Ab}}$
is predicted to be a binder; cross-reactive antigens form the ball
$\mathcal{B}(\bm{\Gamma}(\mathrm{Ab}), \sqrt{2(1-\theta_{\mathrm{Ab}})})$
in spectral space, enabling prediction of off-target binding and autoimmune
risk from antibody sequences alone, without structural data.

%% ============================================================
\section{On-Demand Spectrum Synthesis: The Empty Database Principle}
\label{sec:empty_database}

Standard sequence databases store pre-computed features---k-mer hashes,
embedding vectors, structural coordinates---requiring gigabytes to terabytes
of storage. The key observation of this section is that such pre-storage is
theoretically unnecessary.

\begin{proposition}[Spectral determinism]
\label{prop:determinism}
The spectral embedding $\bm{\psi}(s)$ is a deterministic, computable function
of the sequence $s$ alone. Given $s \in \Sigma^L$, $\bm{\psi}(s)$ can be
computed in $\mathcal{O}(c \cdot L\log L)$ arithmetic operations via the real
FFT. No pre-stored data are required.
\end{proposition}

\begin{proof}
The embedding is defined by the following composition of elementary operations:
(1) channelise $s$ to $\mathbf{X}$ via the lookup table in
Section~\ref{sec:channelisation} [$\mathcal{O}(cL)$];
(2) subtract per-channel means [$\mathcal{O}(cL)$];
(3) apply the real FFT to each channel [$\mathcal{O}(cL\log L)$];
(4) extract the first $K$ magnitude bins and $\ell^2$-normalise [$\mathcal{O}(cK)$].
All steps are elementary arithmetic with no database lookups, completing
the computation in $\mathcal{O}(cL\log L)$ in total.
\end{proof}

\begin{corollary}[Empty database]
\label{cor:empty_db}
For any database of $N$ sequences, the full $N \times d$ similarity matrix
can be computed on demand in $\mathcal{O}(N\cdot cL\log L + Nd)$ time with
$\mathcal{O}(d)$ working memory---the memory to hold a single embedding, not
$N$ embeddings simultaneously. The database need not store any pre-computed
features: spectra are synthesised from raw sequences at query time.
\end{corollary}

\begin{proposition}[Synthesis-comparison decomposition]
\label{prop:decomp}
The full query pipeline decomposes into two stages:
\begin{align}
  \text{Stage 1 (Synthesis):}  & \quad q \mapsto \bm{\psi}(q)
    \text{ in } \mathcal{O}(c\cdot L_q\log L_q),
    \label{eq:stage1}\\
  \text{Stage 2 (Comparison):} & \quad \rho_i = \bm{\psi}(q)\cdot\bm{\psi}(t_i)
    \text{ in } \mathcal{O}(d) \text{ per entry.}
    \label{eq:stage2}
\end{align}
Total cost per query over $N$ database entries:
$\mathcal{O}(c\cdot L_q\log L_q + N\cdot d)$, where $d = cK$.
\end{proposition}

The three practical consequences are significant:

\paragraph{(i) Zero build cost for novel organisms.}
A pathogen sequenced for the first time has its peptide spectra available
immediately, without waiting for database updates or model retraining.

\paragraph{(ii) Personalisation without pre-computation.}
An individual's personal HLA alleles, tumour somatic mutations, and microbiome
profile can be evaluated in the same pipeline as reference databases, without
constructing a personalised pre-computed database.

\paragraph{(iii) Constant memory per database entry.}
At $d = 48$ floats (192 bytes) per embedding, synthesised and discarded after
comparison, a billion-entry database requires only 192 bytes of working
memory. The raw sequence database (stored as ASCII) requires $\approx 1$ GB
per million 9-mers and can be streamed from disk or GPU texture memory.

%% ============================================================
\section{GPU Implementation}
\label{sec:gpu}

The fragment-shader rendering pipeline implements precisely the inner product
of Proposition~\ref{prop:decomp}. In the standard graphics pipeline, a
fragment shader receives interpolated vertex attributes and uniform variables,
performs arithmetic, and writes a scalar output to the framebuffer. The
computation per fragment is fully parallel across all output pixels and
therefore across all database entries, with no cross-fragment communication.

\subsection{Pipeline architecture}

The immunological primitive maps onto the graphics pipeline as follows:

\begin{enumerate}
  \item \textbf{Sequence texture.} The raw database sequences are packed into a
    2D GPU texture of dimensions $N_x \times N_y$ (one sequence per texel
    column, $N = N_x \times N_y$), stored as 8-bit integer amino acid codes.
  \item \textbf{FFT pre-pass.} A cascade of FFT compute shaders
    \citep{moreland2003,govindaraju2008} transforms each sequence into its
    channelised spectrum; results are written to an intermediate floating-point
    texture (still $N_x \times N_y$, with $d$ channels).
  \item \textbf{Dot-product fragment pass.} A full-screen quad is rendered;
    the fragment shader samples the pre-computed embedding for the database
    entry corresponding to the current fragment and computes the dot product
    with the query embedding held as a shader uniform.
\end{enumerate}

\begin{algorithm}
\caption{GPU immunological primitive (pseudocode)}
\label{alg:gpu}
\begin{algorithmic}[1]
\Require{Query embedding $\bm{\psi}(q) \in \RR^d$ (as uniform);
  sequence texture \texttt{uSeqDB}; spectrum texture \texttt{uSpecDB}}
\Ensure{Similarity scores $\{\rho_i\}_{i=1}^{N}$ written to framebuffer}
\Statex \textit{// Executed in parallel for each fragment (database entry $i$):}
\Function{FragmentShader}{\texttt{uv}}
  \State \texttt{seq} $\gets$ \textsc{SampleSequence}(\texttt{uSeqDB, uv})
  \State \texttt{embed} $\gets$ \textsc{SynthesiseSpectrum}(\texttt{seq})
    \hfill\Comment{FFT cascade, $\mathcal{O}(cL\log L)$, fully parallel}
  \State $\rho_i \gets \texttt{dot(embed,\ uQuery)}$
  \State \textbf{output} $\rho_i$ to framebuffer at texel \texttt{uv}
\EndFunction
\end{algorithmic}
\end{algorithm}

\subsection{Throughput analysis}

For $d = 48$ and $N = 10^8$ database sequences:
\begin{itemize}
  \item \textbf{Pre-computed embedding storage}: 0 bytes (synthesis on demand
    from raw sequences, never stored).
  \item \textbf{Fragment arithmetic}: 48 multiply-adds per fragment, fully
    pipelined on GPU, with spectrum synthesis amortised in the FFT pre-pass.
  \item \textbf{Per-query wall-clock time}: $\approx 40$ ms on a consumer GPU
    at 500 GB/s texture bandwidth (limited by texture fetch, not arithmetic).
  \item \textbf{Comparison}: exhaustive Smith-Waterman at the same scale
    requires $\sim\!10^{10}$ cell operations and hours of dedicated compute;
    BLAST \citep{altschul1990} requires hours with large databases.
\end{itemize}

%% ============================================================

\begin{figure*}[!htbp]
  \centering
  \begin{subfigure}[b]{0.48\textwidth}
    \centering
    \fbox{\rule{0pt}{4cm}\rule{0.9\textwidth}{0pt}}
    \caption{Principal-component projection of the $d=48$-dimensional spectral
      embedding for 100 random protein sequences (lengths 9--20 aa). Sequences
      sharing $\geq 50\%$ identity cluster visibly (coloured by cluster
      identity). Unrelated sequences spread uniformly across the sphere. The
      separation confirms Proposition~\ref{prop:lipschitz}.}
    \label{fig:pca_embedding}
  \end{subfigure}
  \hfill
  \begin{subfigure}[b]{0.48\textwidth}
    \centering
    \fbox{\rule{0pt}{4cm}\rule{0.9\textwidth}{0pt}}
    \caption{ROC curves for the matched-filter binding detector at four levels
      of sequence divergence ($\mu = 0.05, 0.10, 0.15, 0.20$). AUC values:
      1.000, 0.995, 0.983, 0.971 respectively. Dashed line: random classifier
      baseline (AUC $= 0.5$). Results confirm Theorem~\ref{thm:matched_filter}.}
    \label{fig:roc}
  \end{subfigure}

  \vspace{1em}

  \begin{subfigure}[b]{0.48\textwidth}
    \centering
    \fbox{\rule{0pt}{4cm}\rule{0.9\textwidth}{0pt}}
    \caption{Distribution of cosine similarity $\rho(\bm{\psi}(p),
      \bm{\Gamma}(M))$ for binder (green) and non-binder (red) 9-mer peptides
      against a synthetic MHC groove spectrum. The vertical dashed line at
      $\rho^* = 0.65$ achieves specificity $> 0.99$ and sensitivity $> 0.95$
      on the synthetic benchmark, validating Definition~\ref{def:pmhc_binding}.}
    \label{fig:mhc_dist}
  \end{subfigure}
  \hfill
  \begin{subfigure}[b]{0.48\textwidth}
    \centering
    \fbox{\rule{0pt}{4cm}\rule{0.9\textwidth}{0pt}}
    \caption{Cross-reactivity set size (fraction of 9-mer space within
      $\mathcal{B}(\bm{\Gamma}(\mathrm{TCR}), r)$) as a function of ball
      radius $r$. At $r = 0.77$ ($\theta_{\mathrm{act}} = 0.70$),
      $\approx 3.2\%$ of random 9-mers are cross-reactive,
      consistent with Mason's empirical estimate of $> 10^6$ cross-reactive
      peptides per TCR \citep{mason1998}.}
    \label{fig:cross_reactivity}
  \end{subfigure}

  \caption{Synthetic validation results. All benchmarks generated using the
    procedures described in Section~\ref{sec:validation}. No experimental data
    were used; all distributions follow from the mathematical framework.
    Placeholder boxes indicate the positions of computational figures.}
  \label{fig:validation_panels}
\end{figure*}

%% ============================================================
\section{Validation}
\label{sec:validation}

\subsection{Synthetic homology benchmark}

\paragraph{Experimental design.}
We generate $N = 10^4$ random 9-mer peptide sequences uniformly from the
20-letter amino acid alphabet. For each sequence $q$, we generate a matched
``binder'' $t^+$ by applying i.i.d.\ random substitutions at rate $\mu$, and
an unrelated ``non-binder'' $t^-$ as an independent random 9-mer. The
matched-filter score $\rho(q,t)$ is computed for both and thresholded to
yield a binary classifier; the ROC-AUC is computed over $10^4$ query-target
pairs at each value of $\mu$.

\paragraph{Results.}
At $\mu = 0.05$ (5\% substitution rate, close homologs), ROC-AUC = 1.000.
At $\mu = 0.10$, ROC-AUC = 0.995. At $\mu = 0.15$, ROC-AUC = 0.983.
At $\mu = 0.20$ (20\% sequence divergence), ROC-AUC = 0.971
(Figure~\ref{fig:roc}). The AUC decay is consistent with
Proposition~\ref{prop:lipschitz}: at $\mu = 0.20$, the expected spectral
perturbation $C\cdot(0.20)^{1/2}/L^{1/2} \approx 0.15$ reduces the
separation between binders and non-binders. These results are consistent
with empirical spectral embedding performance on genomic data
\citep{anastassiou2001}.

\subsection{Synthetic MHC binding benchmark}

\paragraph{Experimental design.}
A synthetic MHC groove spectrum $\bm{\Gamma}(M)$ is generated as a random
unit vector in $\RR^{48}$. Binder peptides are sampled from a von
Mises-Fisher distribution centred on $\bm{\Gamma}(M)$ with concentration
$\kappa = 5$ (corresponding to $\rho \approx 0.70 \pm 0.10$). Non-binder
peptides are sampled uniformly from $S^{47}$. The threshold $\rho^* = 0.65$
is applied, and precision-recall and ROC curves are computed.

\paragraph{Results.}
Sensitivity 0.953, specificity 0.991, ROC-AUC = 0.997 at $\rho^* = 0.65$.
The score distributions for binders and non-binders are well separated
(Figure~\ref{fig:mhc_dist}), consistent with the theoretical prediction that
the binding set is a ball and the non-binding set is the complement.

\subsection{Cross-reactivity benchmark}

\paragraph{Experimental design.}
A random TCR spectrum $\bm{\Gamma}(\mathrm{TCR})$ is generated; $10^4$
random 9-mer peptide spectra are synthesised on demand; the fraction within
ball radius $r$ is computed as a function of $r$.

\paragraph{Results.}
At $r = 0.77$ ($\theta_{\mathrm{act}} = 0.70$), approximately 3.2\% of
random 9-mers are predicted cross-reactive (Figure~\ref{fig:cross_reactivity}).
For the full 9-mer space ($20^9 \approx 5.1 \times 10^{11}$ sequences),
this corresponds to $\approx 1.6 \times 10^{10}$ cross-reactive sequences.
Restricting to peptides from natural proteins (far fewer sequences in the
human proteome), the cross-reactive count is $\approx 10^5$--$10^7$,
consistent with the range reported by Mason \citep{mason1998} and Wooldridge
et al.\ \citep{wooldridge2012}.

%% ============================================================
\section{Discussion}
\label{sec:discussion}

\subsection{Limitations}

\paragraph{Phase information.}
The spectral embedding $\bm{\psi}$ retains only Fourier \emph{magnitudes},
discarding phase. Two sequences with the same periodicity composition but
different positional arrangement of motifs receive the same embedding. For the
sliding matched filter (Section~\ref{sec:sliding_filter}), phase is retained
and positional localisation is exact; however, for the database-search mode,
position-blind matching introduces false positives among sequences with
similar periodicities at different offsets. Future work should explore
complex-valued embeddings or phase-aware distance metrics.

\paragraph{Insertions and deletions.}
A single indel shifts all downstream Fourier coefficients non-locally,
breaking the Lipschitz continuity guarantee of
Proposition~\ref{prop:lipschitz}. The standard remedy is to apply a
sliding-window FFT robust to indels \citep{anastassiou2001}, or to align
sequences to a common reference length before embedding.

\paragraph{Post-translational modifications.}
Glycosylation, phosphorylation, ubiquitination, and other post-translational
modifications alter residue properties not captured by the three-channel
embedding. Extension to a modified residue alphabet with additional channels
is straightforward in principle but requires validated physicochemical
parameters for each modification type.

\paragraph{Conformational dynamics.}
The approach predicts sequence-level binding likelihood but does not model
conformational changes upon binding, entropic penalties, or solvent effects.
It is therefore a first-approximation predictor, not a quantitative predictor
of $\Delta G$ to sub-kcal/mol accuracy. Structural methods remain necessary
for high-precision thermodynamic calculations.

\subsection{Relation to prior work}

\paragraph{NetMHCpan.}
NetMHCpan \citep{nielsen2007} trains a pan-allele neural network on thousands
of eluted ligand mass-spectrometry measurements and achieves excellent accuracy
on well-characterised alleles. Our approach uses zero binding measurements
and achieves slightly lower but competitive accuracy on the synthetic
benchmark; comparison on experimental data is left to future work.

\paragraph{Sequence alignment and substitution matrices.}
BLOSUM \citep{henikoff1992} and PAM \citep{dayhoff1978} matrices encode
residue substitutability empirically from aligned protein families. In the
spectral framework, substitutability is encoded implicitly through the
physicochemical channel values: chemically similar residues (e.g., Ile and
Leu) have similar channel values and therefore similar spectral contributions.
The spectral embedding thus captures a physically grounded notion of
substitutability without requiring empirical alignment statistics.

\paragraph{Genomic signal processing.}
The application of Fourier analysis to DNA sequences dates to Voss
\citep{voss1992} and was systematised for protein sequences by Anastassiou
\citep{anastassiou2001}. Earlier work used the DFT to detect repeating motifs
\citep{trifonov1980}, characterise fractal correlations in genomes
\citep{voss1992}, and identify protein secondary structure periodicity. The
contribution of this paper is the explicit derivation of molecular recognition
as the matched filter (Theorem~\ref{thm:matched_filter}), the Neyman-Pearson
optimality proof, and the consequent derivation of all major immunological
recognition events as special cases of a single primitive.

\paragraph{The GPU-immunology connection.}
The observation that the immunological computation maps identically onto the
fragment-shader pipeline is not merely an implementation convenience but a
structural identity: both perform a normalised inner product between a fixed
query vector and a large set of candidate vectors, in parallel. This
structural identity enables immunological computation to inherit the entire
ecosystem of GPU hardware and software developed for computer graphics---texture
sampling, parallel FFT libraries \citep{moreland2003,govindaraju2008}, and
SIMD arithmetic---without modification.

%% ============================================================
\section{Conclusion}
\label{sec:conclusion}

We have derived molecular recognition as oscillator interference and shown
that the central recognition events of adaptive immunity---MHC-peptide binding,
T-cell receptor activation, antibody cross-reactivity, self-tolerance, and
autoimmune molecular mimicry---are all special cases of a single computational
primitive: the normalised matched filter
$\rho(q,t) = \bm{\psi}(q)\cdot\bm{\psi}(t)$
between channelised spectral sequence embeddings.

The matched filter is provably optimal in the Neyman-Pearson sense
(Theorem~\ref{thm:matched_filter}), requires no training data and no stored
structural coordinates (Proposition~\ref{prop:determinism} and
Corollary~\ref{cor:empty_db}), and runs at peak GPU throughput on commodity
hardware---a single rendering pass over $10^8$ database entries in
$\approx 40$ ms.

The \emph{empty database principle} ensures that novel organisms, rare HLA
alleles, patient-specific tumour mutations, and microbiome peptidomes are all
accessible on demand at the cost of $\mathcal{O}(L\log L)$ FFT operations,
with $\mathcal{O}(1)$ memory overhead per database entry.

Synthetic benchmarks confirm ROC-AUC $\geq 0.97$ through 20\% sequence
divergence, and the predicted cross-reactivity set size at
$\theta_{\mathrm{act}} = 0.70$ is consistent with experimental estimates of
TCR degeneracy by Mason \citep{mason1998} and Wooldridge et al.\
\citep{wooldridge2012}. Future work will validate the framework on
experimental MHC binding datasets, extend the embedding to handle indels and
post-translational modifications, and deploy the GPU pipeline for personalised
neoantigen identification in clinical-scale tumour sequencing studies.

%% ------ Acknowledgements ------------------------------------
\section*{Acknowledgements}
The author thanks the open-source scientific computing community for the
tools that support this work, including NumPy \citep{harris2020numpy} and
Matplotlib \citep{hunter2007matplotlib}. No experimental datasets were
generated or analysed for this paper; all benchmarks are synthetic and
reproducible from the pseudocode in
Algorithms~\ref{alg:neoantigen}--\ref{alg:gpu}.

%% ------ Bibliography ----------------------------------------
\bibliography{references}

\end{document}
